\section{The modified Ricci concentration regime}\label{sec:modified-ricci-concentration}

Suppose \(g_i\) are unit-volume positive Yamabe metrics on possibly
varying nonstandard smooth homotopy four-spheres \(M_i\), with
\(e_i=E_+(g_i)=E_-(g_i)\le2T-\delta\) for fixed \(\delta>0\).
Use \(\sigma_i=\int L_{g_i}^2dV_{g_i}\) as in
Section~\ref{sec:isotropic-gap}. Consider the residual regime
\(36\sigma_i-Y_i^2\to0\) and the actual modified Yamabe metrics
\(q_i\) of volume one in these classes. The following statement
identifies which shrinking comparisons can carry Ricci energy.

\begin{theorem}[Modified shrinking limits]\label{thm:modified-shrinking-limits}
Let \(a_i\to0\). If \(a_i^{-2}q_i\) converge locally uniformly and
strongly in \(W^{1,4}\) to a complete connected smooth Riemannian
four-manifold under oriented smooth embeddings of neighborhoods
of a compact exhaustion, then the limit is Euclidean. The full
curvature integrals on the images of fixed compacts tend to zero.

The same conclusion holds for the complete pointed harmonic
limit obtained if the curvature of \(a_i^{-2}q_i\) is uniformly
bounded on each fixed-radius ball at the chosen centers, with
bounds depending on the radius. In particular, for
\(K_i=\max_{M_i}|\Rm_{q_i}|_{\rm tens}\), one has \(K_i\to\infty\)
and for each fixed \(A>0\),
\begin{equation}\label{eq:modified-peak-energy}
 \sup_{x\in M_i}\int_{B_{q_i}(x,A/\sqrt{K_i})}
              |\Rm_{q_i}|_{\rm tens}^2dV_{q_i}\longrightarrow0.
\end{equation}
\end{theorem}

\begin{proof}
Pass to a subsequence with \(e_i\to e\) and \(Y_i\to y>0\), using
the strict band for the positive lower bound. As in the proof of
Theorem~\ref{thm:isotropic-gap}, \(R_{q_i}=6L_{q_i}+c_i\),
\(c_i\to0\), \(L_{q_i}\to y/6\), and \(R_{q_i}\to y\) in
the auxiliary \(L^2\) norms for the last two statements. Since
\(|W^\pm|_{\rm HS}^2\le6L^2\), the Weyl and scalar squared
integrals tend to zero on sets of vanishing \(q_i\)-volume.
Gauss--Bonnet gives
\begin{equation}\label{eq:modified-total-ricci}
 \int |Z_{q_i}|^2dV_{q_i}\longrightarrow4e+y^2/12-8T.
\end{equation}
There is also an inhomogeneous Sobolev inequality with uniform
constant, by the ordinary Yamabe inequality and the scalar
\(L^2\) estimate. These facts persist on smooth volume-one
conformal approximations of \(q_i\), chosen separately at each
index. We may require relative distance errors and pointwise
errors in curvature norms tending to zero, and \(o(1)\) total
variation error for the squared curvature measures.

For a supplied marked comparison \((X,h)\), the unrescaled
volume of a fixed compact image is \(O(a_i^4)\). Hence both
rescaled Weyl blocks and the scalar curvature tend to zero
in local \(L^2\). The all-lowered curvature formula passes
distributionally under the stated first-order convergence,
since its quadratic first derivatives converge strongly in
\(L^2\). The full curvature is uniformly bounded in \(L^2\) by
\eqref{eq:modified-total-ricci} and the other curvature budgets.
Consequently \(R_h=W_h=0\), and by exhaustion
\(d_X:=\int_X|\Ric_h|^2dV_h\le4e+y^2/12-8T\).
The ordinary Yamabe inequality passes to compactly supported
tests on \(X\), since the compact scalar term tends to zero by
H\"older. Thus \(y\|f\|_4^2\le6\|df\|_2^2\). This excludes
compactness and gives \(\operatorname{Vol}_h B_s(x)\ge bs^4\)
at all radii by the dyadic cutoff estimate.

The exact scalar-flat conformally flat metric has harmonic
curvature and finite \(\int|\Rm|_{\rm tens}^2=2d_X\). With the finite homogeneous Sobolev constant already proved,
the harmonic-curvature case of Tian--Viaclovsky gives finitely many
ends, each asymptotically locally Euclidean (ALE) of every order
less than two \cite[Theorem~5.1]{TianViaclovsky2008Volume}.

Let \(k\) be the number of ends, \(m_j\) their link-group orders,
and \(m=\max_j m_j\). A whole compact truncation embeds in an
actual source. Its first and second rational homology vanish by
Mayer--Vietoris since the end links are rational homology spheres.
The relative fundamental class and Poincar\'e--Lefschetz duality
give third Betti number \(k-1\), hence \(\chi(X)=2-k\).
The ALE Gauss--Bonnet identity in the scalar-flat conformally
flat case now gives the exact cost
\begin{equation}\label{eq:modified-ale-cost}
 d_X=4T\left(k-2+\sum_{j=1}^k m_j^{-1}\right)
 \quad\cite[Section~7, (7.4)]{TianViaclovsky2005ALE}.
\end{equation}
The homogeneous Yamabe inequality on each end gives
\(y\le S/\sqrt{m_j}\), where \(S=8\sqrt6\pi\). Cut the Euclidean
radial sharp test to an annulus and dilate its support out the
end at fixed cutoffs. On the quotient the gradient
integral is divided by \(m_j\) and the squared \(L^4\) norm by
\(\sqrt{m_j}\). Remove the cutoffs after the metric limit. The
Ricci upper budget and \eqref{eq:modified-ale-cost} now give
\[
 \delta+T\left(k-2+\sum_j m_j^{-1}\right)\le 2T/m.
\]
For \(k\ge2\) the left side is strictly larger than \(2T/m\).
Thus \(k=1\). The nonnegativity of \(d_X\) in
\eqref{eq:modified-ale-cost} forces the single end group to be
trivial and \(d_X=0\). The metric is flat with asymptotic volume
ratio one, and Bishop--Gromov rigidity makes it Euclidean.

To get compact energy convergence, put the two Levi-Civita
connections on the same oriented Euclidean bundle using metric
square roots. Their difference \(\Theta_i\) tends to zero in
\(L^4\). The affine path of connections has curvature forms
\(\Omega_{i,t}=(1-t)\Omega+t\Omega_i-t(1-t)\Theta_i\wedge\Theta_i\)
uniformly bounded in \(L^2\). Differentiating the Pfaffian along
the path and using covariant Bianchi gives the difference of
Euler forms as the derivative of a transgression involving
\(2\int_0^1\operatorname{Pf}(\Theta_i,\Omega_{i,t})dt\), with
the common Euler normalization. Its \(L^{4/3}\) norm tends to
zero. Thus the Euler densities converge against compact smooth
cutoffs. Since the limit is flat and the scalar and Weyl
terms vanish in \(L^1\), the Euler identity with coefficient
\(-1/2\) on \(|Z|^2\) gives vanishing trace-free Ricci integrals
on fixed compacts. The full curvature then also tends to zero
in \(L^2\) there.

For the bounded-curvature extraction, the original Sobolev
estimate rescales with a zero-order coefficient \(O(a_i^2)\).
It gives uniform lower volumes on fixed small balls by cutoff
iteration. Together with curvature bounds on enlarged balls
this gives local injectivity and harmonic-radius estimates
\cite[Theorem~4.7]{CheegerGromovTaylor1982}%
\cite[Theorem~0.3]{AndersonCheeger1992Harmonic}.
The harmonic coordinate Ricci equation and the curvature bound
give local \(W^{2,p}\) bounds for finite \(p>4\), by first gaining
bounded first derivatives from a large-exponent harmonic-radius
bound. A diagonal extraction on buffered windows gives a
complete marked limit in \(C^{1,\beta}\) and weak \(W^{2,p}\),
with the images covering each fixed smaller source ball
\cite[Theorem~3.1 and proof]{BelegradekKapovitch2000Finiteness}.
The scalar and Weyl limits vanish by the previous distributional
argument. For this limit with bounded curvature the weak
curvature identities give \(\Delta_h\Ric_h=\Rm_h*\Ric_h\).
Local tensor elliptic regularity and the harmonic metric equation
bootstrap to smoothness in the same charts
\cite[Lemma~5.1 and local form of Theorem~2.2(ii)]{AvalosCogoRoyoAbrego2026Bach}.
The resulting smooth exhaustion embeddings retain the first-order
convergence. The preceding limit and energy conclusions apply.

If \(K_i\) were bounded on a subsequence, the fixed-cutoff
Ricci tail would vanish and the tail-only isotropic recognition
argument of Section~\ref{sec:isotropic-gap} would identify the
standard type. Thus \(K_i\to\infty\). At scale \(K_i^{-1/2}\)
the full rescaled curvature is uniformly bounded globally.
Any sequence of centers violating \eqref{eq:modified-peak-energy}
would have a buffered harmonic limit with vanishing curvature
integral on the indicated window, a contradiction. The same
statement holds for the unsmoothed \(q_i\) by the relative
distance and squared-curvature measure bounds above.
\end{proof}

If a sequence of balls of radii \(r_i\to0\) carries a fixed
positive modified trace-free Ricci charge, then
\(r_i\sqrt{K_i}\to\infty\) by \eqref{eq:modified-peak-energy}.
Along that sequence the rescaled full curvature must be unbounded
on some fixed enlargement, since uniform bounds on every such
enlargement would give the second extraction in
Theorem~\ref{thm:modified-shrinking-limits}. This is a restriction
on actual charged balls. The loss of uniform integrability alone
gives no such marked ball or bounded-curvature enlargement.
